\chapter{Static proportions}
\label{chap:static}

Section 2 and Appendix A of \cite{D23}: the rate of the static proportions algorithms, from the
large deviation principle of the empirical means. Throughout, the arm $k$ belongs to the
exponential family generated by a measure $\rho_k$ on $\R$ with the statistic $\mathrm{id}$, not
almost everywhere constant, and $\KL(x, \mu_k)$ is extended to every level $x$ by the Cramér rate
function of the member of mean $\mu_k$.

\begin{lemma}[{\cite[Lemma 11]{D23}}]
  \label{lem:ldp_single_arm}
  \lean{Degenne2023Existence.isLDP_empiricalMean_expFamily}
  \leanok
  \uses{def:ldp, def:good_rate, lem:cramer_expFamily}
  Let $\rho_k \ne 0$ and let $\mu_k$ be a mean of the family. The empirical mean of $n$ independent
  samples of its member of mean $\mu_k$ satisfies the large deviation principle with speed $n$ and
  good rate function $x \mapsto \KL(x, \mu_k)$.
\end{lemma}

The hypothesis $\rho_k \ne 0$ is implicit in the paper, whose families are families of probability
measures. Without it the statement is false: the zero measure has the single mean $0$, whose member
is the zero measure (\texttt{notes/statement-corrections.md}, Lemma 11).

\begin{proof}
  \leanok
  \uses{lem:cramer_expFamily}
  Lemma~\ref{lem:cramer_expFamily}: Cramér's theorem for the member of mean $\mu_k$, whose
  cumulant generating function is finite near $0$ since its natural parameter is in the interior
  of the natural parameters; its rate function is $\KL(\cdot, \mu_k)$.
\end{proof}

\begin{theorem}[{\cite[Theorem 12]{D23}}]
  \label{thm:ldp}
  \lean{Degenne2023Existence.isLDP_empMean_of_isStaticProportions}
  \leanok
  \uses{def:static_proportions, lem:ldp_single_arm, lem:ldp_pi, lem:ldp_window}
  Let $\cA$ be a static proportions family with proportions $\omega \in \Simplex^0$ and $\mu \in
  \cD$ a vector of means of the families. The vector $\hat\mu_T$ of the empirical means after $T$
  rounds satisfies the large deviation principle with speed $T$ and good rate function
  $\lambda \mapsto \sum_k \omega_k \KL(\lambda_k, \mu_k)$: for every set $S \subseteq \R^K$,
  \begin{align*}
    -\inf_{\lambda \in \inter S} \sum_k \omega_k \KL(\lambda_k, \mu_k)
    &\le \liminf_{T \to \infty} \frac{1}{T} \log \Prob_{\mu, \cA}(\hat\mu_T \in S), \\
    \limsup_{T \to \infty} \frac{1}{T} \log \Prob_{\mu, \cA}(\hat\mu_T \in S)
    &\le -\inf_{\lambda \in \cl S} \sum_k \omega_k \KL(\lambda_k, \mu_k).
  \end{align*}
\end{theorem}

\begin{proof}
  \leanok
  \uses{lem:ldp_single_arm, lem:ldp_pi, lem:ldp_window}
  In the stack-of-rewards model, the empirical mean of arm $k$ is the mean of the first $N_{T,k}$
  samples of its stack, and $\abs{N_{T,k} - T \omega_k} \le K$. The member of mean $\mu_k$ of the
  family has a cumulant generating function finite near $0$, and its Cramér rate function is
  $\KL(\cdot, \mu_k)$ (Lemma~\ref{lem:ldp_single_arm}). Lemma~\ref{lem:ldp_window}, through the
  bounds of Lemma~\ref{lem:ldp_pi} on the stacks with numbers of samples in a window of size
  $2K + 1$, gives the principle with the sum of the rates. The recommendation rule plays no role.
\end{proof}

\begin{theorem}[{\cite[Theorem 1]{D23}}]
  \label{thm:oracle_difficulty_sp}
  \lean{Degenne2023Existence.tendsto_errorRate_of_isStaticProportions}
  \leanok
  \uses{asm:answer_nhds, thm:ldp, lem:error_set_subset, lem:rate_continuousOn,
    lem:tendsto_errorRate_iff}
  Suppose that the task satisfies Assumption 1 and has a single correct answer on $\cD$, and that
  the samples of each family lie in the closure of its means. Let $\cA$ be a static proportions
  family with proportions $\omega \in \Simplex^0$ whose recommendation rule answers $i$ at every
  vector of the closure of $\prod_k \cM_k$ that is not in the closure of the alternatives to $i$.
  Then for every $\mu \in \cD$,
  \[ \lim_{T \to \infty} h_{\mu, T}(\cA) = \Bigl( \inf_{\lambda \in \Alt(\mu)} \sum_{k=1}^K
    \omega_k \KL(\lambda_k, \mu_k) \Bigr)^{-1}. \]
\end{theorem}

\begin{proof}
  \leanok
  \uses{lem:oracle_difficulty_sp_measurable, thm:ldp, lem:error_set_subset, lem:rate_continuousOn,
    lem:tendsto_errorRate_iff}
  The set of correct answers of $\mu$ is a singleton, measurable, so
  Lemma~\ref{lem:oracle_difficulty_sp_measurable} applies. Its proof is the paper's. Let
  $i = i^{\star}(\mu)$. The algorithm errs when $\hat\mu_T \in \Alt(\mu)$, an open set by
  Assumption 1, and only when $\hat\mu_T$ is in the closed set $F$ of the vectors of
  $\cl(\prod_k \cM_k)$ at which the recommendation is not $i$; $F \subseteq \cl(\Alt(\mu))$
  (Lemma~\ref{lem:error_set_subset}). The rate function is continuous on $\cl(\prod_k \cM_k)$
  (Lemma~\ref{lem:rate_continuousOn}), so its infima over $\Alt(\mu)$, $\cl(\Alt(\mu))$ and $F
  \cup \Alt(\mu)$ agree. Theorem~\ref{thm:ldp} gives $\frac1T \log p_{\mu, T}(\cA) \to -\inf_{\Alt(\mu)}
  \sum_k \omega_k \KL(\lambda_k, \mu_k)$, which is the claim
  (Lemma~\ref{lem:tendsto_errorRate_iff}).
\end{proof}

The paper allows any extension of $i^{\star}$ (E1 of \texttt{notes/paper-errata.md}) and uses that
the empirical means lie in the closure of the means (E2).

The Lean statements of Theorems~\ref{thm:oracle_difficulty_sp} and~\ref{thm:oracle_difficulty_class}
assume that the answers have measurable singletons, as the paper's answers, a finite set of indices
every subset of which is an event. Without it they are false: for the trivial $\sigma$-algebra on
the answers, the probability of error, an outer measure, is $1$ for every algorithm
(\texttt{notes/statement-corrections.md}, Theorems 1 and 5). The proof of
Theorem~\ref{thm:oracle_difficulty_sp} needs neither this nor finitely many answers, only that the
set of correct answers of the model is measurable (Lemma~\ref{lem:oracle_difficulty_sp_measurable}).

\begin{lemma}[{\cite[Theorem 1]{D23}}]
  \label{lem:oracle_difficulty_sp_measurable}
  \lean{Degenne2023Existence.tendsto_errorRate_of_isStaticProportions_of_measurableSet}
  \leanok
  \uses{asm:answer_nhds, thm:ldp, lem:error_set_subset, lem:error_event, lem:errorProb_static,
    lem:static_error_lower, lem:exists_rec, lem:rate_continuousOn, lem:tendsto_errorRate_iff}
  Theorem~\ref{thm:oracle_difficulty_sp} holds at every model $\mu \in \cD$ whose set of correct
  answers is measurable, without the finiteness of the answers.
\end{lemma}

\begin{proof}
  \leanok
  \uses{thm:ldp, lem:error_set_subset, lem:error_event, lem:errorProb_static,
    lem:static_error_lower, lem:exists_rec, lem:rate_continuousOn, lem:tendsto_errorRate_iff}
  The paper's proof of Theorem~\ref{thm:oracle_difficulty_sp}, with the probability of error as the
  probability that the recommendation of $\hat\mu_T$ is wrong (Lemma~\ref{lem:errorProb_static}).
  Lower bound: the rule is correct on $\cD$ (Lemma~\ref{lem:exists_rec}), and
  Lemma~\ref{lem:static_error_lower} applies. Upper bound: for $T$
  large, the probability of error is at most $\Prob(\hat\mu_T \in F)$
  (Lemma~\ref{lem:error_event}), and the rate function is at least
  $\inf_{\Alt(\mu)} \sum_k \omega_k \KL(\lambda_k, \mu_k)$ on $F$ by its continuity on
  $\cl(\prod_k \cM_k)$ (Lemma~\ref{lem:rate_continuousOn}). The closure and interior conditions of
  Assumption 1 are not used.
\end{proof}

\begin{lemma}[{\cite[proof of Theorem 1]{D23}}]
  \label{lem:error_event}
  \lean{Degenne2023Existence.eventually_errorProb_le_of_isStaticProportions}
  \leanok
  \uses{lem:error_set_subset, lem:errorProb_static, lem:empMean_mem}
  Suppose that the samples of each family lie in the closure of its means, and that the
  recommendation rule is as in Theorem~\ref{thm:oracle_difficulty_sp}. At a model $\mu \in \cD$
  with the single and measurable correct answer $i$, for $T$ large,
  $p_{\mu, T}(\cA) \le \Prob_{\mu, \cA}(\hat\mu_T \in F)$, with
  $F = \cl(\prod_k \cM_k) \cap \cl(\Alt_i)$.
\end{lemma}

\begin{proof}
  \leanok
  \uses{lem:error_set_subset, lem:errorProb_static, lem:empMean_mem}
  For $T > K / \min_k \omega_k$ every arm is pulled, since $N_{T,k} \ge T \omega_k - K$, so the
  empirical means lie in $\cl(\prod_k \cM_k)$ almost surely (Lemma~\ref{lem:empMean_mem}: the
  members of the families are absolutely continuous with respect to $\rho_k$). There, the rule
  errs only on $F$ (Lemma~\ref{lem:error_set_subset}).
\end{proof}

\begin{theorem}[{\cite[Theorem 1]{D23}}]
  \label{thm:oracle_difficulty_class}
  \lean{Degenne2023Existence.oracleDifficulty_staticClass}
  \leanok
  \uses{thm:oracle_difficulty_sp, def:static_class, def:static_difficulty,
    lem:oracleDifficulty_staticClass_le, lem:staticClass_le_oracleDifficulty}
  Under the hypotheses of Theorem~\ref{thm:oracle_difficulty_sp} on the families and the task,
  for every $\mu \in \cD$,
  \[ H_{\Csp}(\mu) = \inf_{\omega \in \Simplex^0} \lim_{T \to \infty} h_{\mu, T}(\cA^{sp}_{\omega})
    = \Bigl( \max_{\omega \in \Simplex} \inf_{\lambda \in \Alt(\mu)} \sum_{k=1}^K \omega_k
    \KL(\lambda_k, \mu_k) \Bigr)^{-1} = \Hsp(\mu). \]
\end{theorem}

\begin{proof}
  \leanok
  \uses{lem:oracleDifficulty_staticClass_le, lem:staticClass_le_oracleDifficulty}
  The set of correct answers of $\mu$ is a singleton, measurable; combine
  Lemmas~\ref{lem:oracleDifficulty_staticClass_le} and~\ref{lem:staticClass_le_oracleDifficulty}.
  The condition of Assumption 1 that each $\cD_i$ is the interior of its closure is not used.
\end{proof}

\begin{lemma}[{\cite[proof of Theorem 1]{D23}}]
  \label{lem:static_error_lower}
  \lean{Degenne2023Existence.neg_le_liminf_log_errorProb_of_isStaticProportions}
  \leanok
  \uses{asm:answer_nhds, thm:ldp, lem:errorProb_static}
  Suppose that the task has open answer sets and a single correct answer on $\cD \subseteq
  \prod_k \cM_k$. Let $\cA$ be a static proportions family with proportions $\omega \in
  \Simplex^0$ whose recommendation rule is correct on $\cD$. At every $\mu \in \cD$ whose set of
  correct answers is measurable,
  \[ -\inf_{\lambda \in \Alt(\mu)} \sum_k \omega_k \KL(\lambda_k, \mu_k) \le \liminf_{T \to \infty}
    \frac1T \log p_{\mu, T}(\cA). \]
\end{lemma}

\begin{proof}
  \leanok
  \uses{thm:ldp, lem:errorProb_static}
  $\Alt(\mu)$ is the union of the open sets $\cD_j$, $j \ne i^{\star}(\mu)$, and at a point of
  $\cD_j$ the rule answers $j$, which is not correct for $\mu$. The lower bound of
  Theorem~\ref{thm:ldp} on this open set, on which the rate function is the cost, gives the claim.
\end{proof}

\begin{lemma}[{\cite[proof of Theorem 1]{D23}}]
  \label{lem:staticClass_le_oracleDifficulty}
  \lean{Degenne2023Existence.staticDifficulty_le_oracleDifficulty_staticClass}
  \leanok
  \uses{lem:static_error_lower, lem:le_liminf_errorRate, def:static_class, def:static_difficulty}
  Under the hypotheses of Lemma~\ref{lem:static_error_lower} on the families and the task,
  $\Hsp(\mu) \le H_{\Csp}(\mu)$ at every $\mu \in \cD$ whose set of correct answers is measurable.
\end{lemma}

\begin{proof}
  \leanok
  \uses{lem:static_error_lower, lem:le_liminf_errorRate}
  A family of $\Csp$ has proportions $\omega \in \Simplex$ and a rule correct on $\cD$, so
  Lemmas~\ref{lem:static_error_lower} and~\ref{lem:le_liminf_errorRate} bound its rate from below by
  $(\inf_{\lambda \in \Alt(\mu)} \sum_k \omega_k \KL(\lambda_k, \mu_k))^{-1} \ge \Hsp(\mu)$.
\end{proof}

\begin{lemma}[{\cite[proof of Theorem 1]{D23}}]
  \label{lem:oracleDifficulty_staticClass_le}
  \lean{Degenne2023Existence.oracleDifficulty_staticClass_le}
  \leanok
  \uses{lem:oracle_difficulty_sp_measurable, lem:exists_static, lem:exists_rec,
    lem:iSup_open_simplex}
  Under the hypotheses of Theorem~\ref{thm:oracle_difficulty_sp} on the families and the task,
  with finitely many answers, $H_{\Csp}(\mu) \le \Hsp(\mu)$ at every $\mu \in \cD$ whose set of
  correct answers is measurable.
\end{lemma}

\begin{proof}
  \leanok
  \uses{lem:oracle_difficulty_sp_measurable, lem:exists_static, lem:exists_rec,
    lem:iSup_open_simplex}
  For every $\omega \in \Simplex^0$, the static proportions family of
  Lemma~\ref{lem:exists_static} with the rule of Lemma~\ref{lem:exists_rec} is in $\Csp$, and its
  rate tends to $(\inf_{\lambda \in \Alt(\mu)} \sum_k \omega_k \KL(\lambda_k, \mu_k))^{-1}$
  (Lemma~\ref{lem:oracle_difficulty_sp_measurable}). The supremum over $\Simplex^0$ equals the
  maximum over $\Simplex$ (Lemma~\ref{lem:iSup_open_simplex}).
\end{proof}

\begin{lemma}
  \label{lem:error_set_subset}
  \uses{asm:answer_nhds}
  Under Assumption 1 with finitely many answers and a single correct answer on $\cD$, let the
  recommendation rule answer $i$ at every vector of $\cl(\prod_k \cM_k)$ outside
  $\cl(\Alt_i)$. Then the vectors of $\cl(\prod_k \cM_k)$ at which it does not answer $i$ lie in
  $\cl(\Alt_i)$.
\end{lemma}

\begin{proof}
  The contrapositive of the hypothesis.
\end{proof}

\begin{lemma}[{\cite[proof of Theorem 1]{D23}}]
  \label{lem:exists_rec}
  \lean{Degenne2023Existence.exists_measurable_rec,
    Degenne2023Existence.eq_of_notMem_closure_alt, Degenne2023Existence.rec_mem_correct}
  \leanok
  \uses{asm:answer_nhds}
  Under Assumption 1 with finitely many answers and a single correct answer on $\cD$, there is a
  measurable recommendation rule that answers $i$ at every vector of $\cl(\prod_k \cM_k)$ outside
  $\cl(\Alt_i)$. Every such rule gives the correct answer on $\cD$.
\end{lemma}

\begin{proof}
  \leanok
  A vector $x$ of $\cl(\prod_k \cM_k) \subseteq \bigcup_k \cl(\cD_k)$ (a finite union of closed
  sets) lies in some $\cl(\cD_k)$, and $\cD_k \subseteq \Alt_i$ for every $i \ne k$, so $x$ is
  outside $\cl(\Alt_i)$ for at most one answer $i$. Answer that $i$ when there is one, and a fixed
  answer otherwise: the set of the vectors with answer $i$ is a Boolean combination of the closed
  sets $\cl(\prod_k \cM_k)$ and $\cl(\Alt_j)$, so the rule is measurable. A model of $\cD$ with
  answer $i$ lies in the open set $\cD_i$, which is disjoint from $\Alt_i$, hence from its
  closure.
\end{proof}
